\documentclass[10pt,a4paper]{amsart}
\usepackage[margin=19mm]{geometry}
\usepackage{amsmath,amssymb,mathtools}
\usepackage[scr=rsfs,cal=euler]{mathalfa}
\usepackage{xfrac}
\usepackage[unicode,hidelinks,bookmarks=false]{hyperref}
\newcommand{\ie}{i.e.\ }
\newcommand{\cf}{cf.\ }
\newcommand{\resp}{respectively\ }
\newcommand{\Subsection}{\subsection}
\providecommand{\nobreakdash}{\nobreak\hbox{-}}
\setlength{\emergencystretch}{2em}
\multlinegap0pt
\usepackage{bm}
\usepackage{bbm}
\usepackage{dsfont}
\usepackage{xspace}
\usepackage{cancel}
\usepackage{mathtools}
\usepackage{amsthm}
\makeatletter
\@ifundefined{proposition}{\newtheorem{proposition}{Proposition}}{}
\makeatother
\newcommand{\E}{\mathbb{E}}
\newcommand{\Gscr}{\mathscr{G}}
\newcommand{\PG}{\mathrm{PG}}
\title[Weak arcs and applications to the DNA-based storage access p]{Weak arcs and applications to the DNA-based storage access problem}
\author{Geertrui Van de Voorde, Ferdinando Zullo}
\date{Source: arXiv:2608.19550v1 (CC BY 4.0)}
\begin{document}
\maketitle

\noindent\emph{Source and licence.} Weak arcs and applications to the DNA-based storage access problem. Version arXiv:2608.19550v1 (CC BY 4.0), distributed under the Creative Commons Attribution 4.0 International licence, as recorded in the arXiv licence metadata for this exact version and shown on the abstract page https://arxiv.org/abs/2608.19550v1. Full rights notice: licenses/arxiv-2608.19550-CC-BY-4.0.txt.\par\smallskip

\noindent\textit{Editorial scope.} Counted unit: the Proposition whose source label is \texttt{lem:weakarc} in the article (source0.tex lines 1645--1668, source label \texttt{lem:weakarc}), delivered together with the article's own complete original proof of it (source0.tex lines 1670--1903, one \texttt{proof} environment of 234 lines). No mathematics is written, repaired or re-derived: statement and proof are the article's own text line for line. Required context retained uncounted: eq:beta (lines 906--908); eq:Tdef (lines 954--956). Every definition, display and label of those lines is retained unchanged. Omitted from this package and not redistributed: 1--905, 909--953, 957--1644, 1669--1669, 1904--2941. Nothing inside a retained excerpt is omitted or altered: every retained body is delivered exactly as the article's lines are written, so no omission receipt of any kind accompanies this package. Citations: the retained excerpts carry no citation command at all, so no bibliography entry of the article is transcribed and no reference is redistributed. Reference adaptation: none was needed and none was made; every reference inside the delivered unit resolves inside it, so no unresolved reference or citation is produced. Printed slips kept literal: no printed word, symbol, hypothesis or number of the counted statement or of its proof was altered. Layout and class: the article's own class is \texttt{article} with packages including graphicx, geometry, booktabs, float, mathrsfs, amsthm, amssymb, amsmath, xcolor, tabularx, array, hyperref, microtype; the wrapper is the lane's pinned \texttt{amsart} editorial wrapper at 10pt on A4, which restates the theorem-like environments the retained text needs and adds the article's own macro definitions that the retained text needs (\texttt{\textbackslash E}, \texttt{\textbackslash Gscr}, \texttt{\textbackslash PG}), with extra wrapper packages (bm, bbm, dsfont, xspace, cancel, mathtools). The delivered numbering is the unit's own. Assets: no retained span contains a figure environment, a graphics-inclusion command, a PDF-page inclusion, a table or a TikZ picture; no image file is copied into this package, the package declares no asset dependency and no image file is redistributed with it. Licence: version arXiv:2608.19550v1 of this article is distributed under the Creative Commons Attribution 4.0 International licence, as recorded in the arXiv licence metadata for this exact version and shown on the abstract page https://arxiv.org/abs/2608.19550v1. A full rights notice is delivered at licenses/arxiv-2608.19550-CC-BY-4.0.txt. Characters: every retained excerpt is pure ASCII, so the delivered engine, XeLaTeX with its default Latin Modern text font, reports no missing character.
\begin{equation}\label{eq:beta}
\E[\tau_{P_i}(\Gscr)] \;=\; 1 \;+\; \sum_{s=1}^{n-1}\frac{\beta_{P_i}(s)}{\binom{n-1}{s}},
\end{equation}

\begin{align}
T(m):=\sum_{s=3}^{n-1}\frac{\binom{m}{s}}{\binom{n-1}{s}}. \label{eq:Tdef}
\end{align}

\begin{proposition}\label{lem:weakarc}
Let $\Gscr\subseteq\PG(3,q)$ be as above then for each fundamental point $P_i$, then
\[
\E[\tau_{P_i}(\Gscr)]
=\frac{13713524a^{6}-13855940a^{5}+5676627a^{4}-1207089a^{3}+140575a^{2}-8515a+210}
{14(5a-1)(7a-1)(9a-1)(9a-2)(10a-1)(10a-3)} ,
\]
% or equivalently, in partial fractions,
% \[
% \E[\tau_F(\Gscr)]
% =\frac{3428381}{992250}
% -\frac{48}{125(5a-1)}
% +\frac{30}{49(7a-1)}
% +\frac{80}{81(9a-2)}
% -\frac{20}{81(9a-1)}
% +\frac{162}{125(10a-3)}
% +\frac{6}{125(10a-1)} .
% \]
which only depends on $a$. It decreases monotonically from
$\tfrac{79}{21}\approx3.7619$ at $a=1$ to
\[
\lim_{a\to\infty}\E[\tau_{P_i}(\Gscr)]=\frac{3428381}{992250}\approx3.4551585 ;\qquad \lim_{a\to\infty}\E[\tau_{P_i}(\Gscr)]/4\approx 0.863788.
\]
\end{proposition}

\begin{proof}
By the symmetry of $\Gscr$ it is enough to compute $\E[\tau_{P_1}]$, using \eqref{eq:beta}. Since a
subset containing any copy of $P_1$ already spans $P_1$, only the $9a$ points distinct from $P_1$
contribute to $\beta_{P_1}$; these are the three vertices $P_2,P_3,P_4$ (each of multiplicity $a$)
and the $6a$ non-fundamental points.

The subspaces not containing $P_1$, spanned by subsets of $\Gscr$ are as folllows. The intersection sizes mentioned are counted with multiplicity.

\begin{itemize}
    \item Points: the three vertices $P_2,P_3,P_4$, each having intersection size $a$ with $\Gscr.$
    \item Lines: the three opposite edges $P_2P_3,P_2P_4,P_3P_4$ each with $2a+a=3a$ points of $\Gscr$; and the $9a$ lines joining a fundamental point
    $P_j$ $(j\ne1)$ to a non-fundamental point of one of the three edges not through $P_j$, each meeting $\Gscr$ in $a+1$ points. 
    \item Planes: the opposite face $P_2P_3P_4$ with $6a$ points of $\Gscr$, the $3a$ planes through one of the edges $P_2P_3,P_2P_4,P_3P_4$ together with one further non-fundamental point (each with $3a+1$ points of $\Gscr$), the $9a^2$ planes through one fundamental point and two non-fundamental points of $\Gscr$, each with $a+2$ points of $\Gscr$, and finally the $16a^3$ planes with $3$ non-fundamental points.
\end{itemize}


Using this, we see that $\beta_{P_1}(\Gscr,0)=1$, $\beta_{P_1}(\Gscr,1)=9a$. For $s=2$, all pairs of points taken from the $9a$ points different from $P_1$ will span a subspace not through $P_1$, except those contained in $P_1P_2,P_1P_3$ or $P_1P_4$. Of the $\binom{2a}{2}$ choices of pairs on $P_1P_i$, we need to exclude those made up of twice $P_i$, so we need to subtract $\binom{a}{2}$:
\[
\beta_{P_1}(\Gscr,2)
=\binom{9a}{2}
-3\left[\binom{2a}{2}-\binom{a}{2}\right]
=3a(12a-1).
\]
For $s\geq 3$, we count the subsets of size $s$ which do not contain $P_1$
in their span according to the dimension of the subspace they span. The subsets spanning a point contribute $3\binom{a}{s}$.
The subsets spanning a line are as follows. The three opposite edges
$P_2P_3,P_2P_4,P_3P_4$ each contain $3a$ points of $\Gscr$, counted with
multiplicity. On each such edge we must subtract the subsets consisting
entirely of copies of one of its two fundamental points. This gives
\[
3\left(\binom{3a}{s}-2\binom{a}{s}\right).
\]
Similarly, the $9a$ lines joining one of the fundamental points
$P_2,P_3,P_4$ to a chosen non-fundamental point on one of the three sides
not through that fundamental point contribute
\[
9a\left(\binom{a+1}{s}-\binom{a}{s}\right).
\]

We now count the subsets spanning a plane. The opposite face $P_2P_3P_4$
contains $6a$ points of $\Gscr$, counted with multiplicity. The subsets
contained in this face but not spanning the face are the subsets spanning a point of
 $P_2,P_3,P_4$, the subsets spanning one of $P_2P_3,P_3P_4,P_2P_4$,
and the subsets spanning one of the $3a$ lines joining one of $P_2,P_3,P_4$ to a chosen
non-fundamental point on the opposite side of the face $P_2P_3P_4$. Hence, $P_2P_3P_4$
contributes
\[
\binom{6a}{s}
-3\binom{a}{s}
-3\left(\binom{3a}{s}-2\binom{a}{s}\right)
-3a\left(\binom{a+1}{s}-\binom{a}{s}\right).
\]

There are $3a$ planes through one of the edges
$P_2P_3,P_2P_4,P_3P_4$ together with one chosen non-fundamental point on the
opposite side through $P_1$. Each such plane contains $3a+1$ points of
$\Gscr$. In such a plane, the proper subspaces spanned by subsets of those $3a+1$ points are the subsets contained in the edge $P_iP_j$ itself, the two lines joining the chosen non-fundamental point to one of the two fundamental points $P_i,P_j$, and the subsets composed of $s$ copies of either $P_i$ or $P_j$. Thus these planes contribute
\[
3a\left(
\binom{3a+1}{s}
-\left(\binom{3a}{s}-2\binom{a}{s}\right)
-2\left(\binom{a+1}{s}-\binom{a}{s}\right)-2\binom{a}{s}
\right).
\]


There are $9a^2$ planes through one fundamental point and two chosen
non-fundamental points of $\Gscr$. Each contains $a+2$ points of $\Gscr$.
These
planes contribute
\[
9a^2\left(
\binom{a+2}{s}
-2\left(\binom{a+1}{s}-\binom{a}{s}\right)-\binom{a}{s}
\right).
\]


Finally, the $16a^3$ planes with three non-fundamental points contribute
\[
16a^3\binom{3}{s}.
\]

Therefore, for $s\geq 3$,
\[
\begin{aligned}
\beta_{P_1}(\Gscr,s)
={}&
3\binom{a}{s}
+3\left(\binom{3a}{s}-2\binom{a}{s}\right)
+9a\left(\binom{a+1}{s}-\binom{a}{s}\right)\\
&{}+\binom{6a}{s}
-3\binom{a}{s}
-3\left(\binom{3a}{s}-2\binom{a}{s}\right)
-3a\left(\binom{a+1}{s}-\binom{a}{s}\right)\\
&{}+3a\left(
\binom{3a+1}{s}
-\binom{3a}{s}
-2\binom{a+1}{s}
+2\binom{a}{s}
\right)\\
&{}+9a^2\left(
\binom{a+2}{s}
-2\binom{a+1}{s}
+\binom{a}{s}
\right)
+16a^3\binom{3}{s}.
\end{aligned}
\]


\[
\begin{aligned}
\beta_{P_1}(\Gscr,s)
={}&
\binom{6a}{s}
+3a\binom{3a+1}{s}
+9a^2\binom{a+2}{s}
+16a^3\binom{3}{s}\\
&{}-3a\binom{3a}{s}
-18a^2\binom{a+1}{s}
+9a^2\binom{a}{s}.
\end{aligned}
\]
We find, using \eqref{eq:Tdef}

\[
\begin{aligned}
\E[\tau_{P_1}(\Gscr)]
={}&1+\frac{9a}{10a-1}
+\frac{3a(12a-1)}{(10a-1)(5a-1)}\\
&{}+T(6a)
+3aT(3a+1)
+9a^2T(a+2)
+16a^3T(3)\\
&{}-3aT(3a)
-18a^2T(a+1)
+9a^2T(a).
\end{aligned}
\]

We find the following values, with $n=10a$,
% \[
% T(6a)
% =
% \frac{3(6a-1)(3a-1)}
% {2(10a-1)(5a-1)},\qquad
% T(3a+1)
% =
% \frac{3a(3a-1)(3a+1)}
% {2(10a-1)(5a-1)(7a-1)}\]

% \[
% T(a+2)
% =
% \frac{a(a+1)(a+2)}
% {2(10a-1)(5a-1)(9a-2)},\qquad
% T(3)
% =
% \frac{3}
% {(10a-1)(5a-1)(10a-3)}.
% \]
% \[
% T(3a)
% =
% \frac{3(3a-1)(3a-2)}
% {14(10a-1)(5a-1)}.
% \]
% \[
% T(a+1)
% =
% \frac{a(a-1)(a+1)}
% {2(10a-1)(5a-1)(9a-1)}.
% \]
% \[
% T(a)
% =
% \frac{(a-1)(a-2)}
% {18(10a-1)(5a-1)}.
% \]

\[
\begin{array}{ll}
T(6a)=\dfrac{3(6a-1)(3a-1)}{2(10a-1)(5a-1)},
   & \quad
T(3a+1)=\dfrac{3a(3a-1)(3a+1)}{2(10a-1)(5a-1)(7a-1)},\\[2ex]
T(a+2)=\dfrac{a(a+1)(a+2)}{2(10a-1)(5a-1)(9a-2)},
   & \quad
T(3)=\dfrac{3}{(10a-1)(5a-1)(10a-3)},\\[2ex]
T(3a)=\dfrac{3(3a-1)(3a-2)}{14(10a-1)(5a-1)},
   & \quad
T(a+1)=\dfrac{a(a-1)(a+1)}{2(10a-1)(5a-1)(9a-1)},\\[2ex]
T(a)=\dfrac{(a-1)(a-2)}{18(10a-1)(5a-1)}.
   &
\end{array}
\]
Hence, 
\[
\begin{aligned}
\E[\tau_{P_1}(\Gscr)]
={}&1+\frac{9a}{10a-1}
+\frac{3a(12a-1)}{(10a-1)(5a-1)}
+\frac{3(6a-1)(3a-1)}
{2(10a-1)(5a-1)}+\frac{9a^2(3a-1)(3a+1)}
{2(10a-1)(5a-1)(7a-1)}\\
&{}+\frac{9a^3(a+1)(a+2)}
{2(10a-1)(5a-1)(9a-2)}+\frac{48a^3}
{(10a-1)(5a-1)(10a-3)}-\frac{9a(3a-1)(3a-2)}
{14(10a-1)(5a-1)}\\
&{}-\frac{9a^3(a-1)(a+1)}
{(10a-1)(5a-1)(9a-1)}+\frac{a^2(a-1)(a-2)}
{2(10a-1)(5a-1)}.
\end{aligned}
\]

In partial fractions,
\[
\begin{aligned}
\E[\tau_{P_1}(\Gscr)]
={}&
\frac{3428381}{992250}
+\frac{6}{125(10a-1)}
+\frac{162}{125(10a-3)}\\
&{}-\frac{20}{81(9a-1)}
+\frac{80}{81(9a-2)}
+\frac{30}{49(7a-1)}
-\frac{48}{125(5a-1)}.
\end{aligned}
\]


Letting $a\to\infty$ gives
$\tfrac{3428381}{992250}\approx3.4551585$, to which the expectation decreases .
\end{proof}

\end{document}
